%=======================================================================
%  Literature Review
%=======================================================================
\section{Literature Review}
\label{sec:litreview}

\subsection{Scope and method of the review}

This review covers four bodies of work that the dissertation draws
together, and it is organised by the question each body answers rather
than chronologically.

\begin{enumerate}[leftmargin=1.6em,itemsep=2pt]
  \item \textbf{What has been built, and what is claimed for it?}
        Electronic and blockchain-based voting systems, from
        cryptographic verifiable-voting schemes to ledger-backed
        proposals (Sections~\ref{sec:lit-evoting}--\ref{sec:lit-bcvoting}).
  \item \textbf{What breaks when it is deployed?} The empirical
        security literature: analyses of fielded internet-voting systems
        and the critical response to blockchain voting specifically
        (Section~\ref{sec:lit-critique}).
  \item \textbf{How is manipulation modelled formally?} Bribery as a
        computational problem, coercion and receipt-freeness, Sybil
        identity attacks, and the statistical forensics of fraudulent
        tallies (Sections~\ref{sec:lit-bribery}--\ref{sec:lit-forensics}).
  \item \textbf{What tools exist for reasoning about a rational,
        resource-bounded attacker?} Concentration inequalities and
        game-theoretic security analysis
        (Section~\ref{sec:lit-probgame}).
\end{enumerate}

Sources were located by targeted search of the primary security venues
(USENIX Security, ACM CCS, Financial Cryptography, IEEE S\&P and their
workshops), the artificial-intelligence and economics literature for
the manipulation-theory strand, and two recent systematic reviews of
blockchain voting~\cite{ohize2025,rahul2025} used as entry points into
the applied literature. Priority was given to peer-reviewed primary
sources; where a widely-cited result exists only as a technical report
or research blog post, as with the Dark DAO
construction~\cite{daian2018}, this is stated explicitly rather than
disguised by citation formatting.

Section~\ref{sec:lit-synthesis} closes the review by identifying what
the four bodies of work leave unaddressed, which is the gap this
dissertation occupies.

\subsection{Electronic voting before blockchain}
\label{sec:lit-evoting}

The security requirements for an electronic election were substantially
settled before distributed ledgers entered the conversation, and it is
worth restating them because blockchain proposals are frequently
evaluated against a narrower list.

\textbf{Verifiability.} A voter should be able to confirm that their
ballot was recorded as cast, and any observer should be able to confirm
that the recorded ballots were tallied correctly. Adida's
Helios~\cite{adida2008} is the canonical open-audit implementation: it
publishes encrypted ballots to a public bulletin board and proves the
tally correct with homomorphic encryption and zero-knowledge proofs, so
that correctness rests on mathematics rather than on trusting the
election authority. Helios matters to the present work as a
counterexample to a common framing: it delivers public verifiability
with no ledger, no chain and no consensus mechanism, which shows that
the verifiability argument for blockchain voting is not an argument for
blockchain specifically.

\textbf{Ballot secrecy and its stronger forms.} Secrecy alone is
insufficient if a voter can \emph{prove} how they voted, because a
provable vote is a sellable vote. Benaloh and
Tuinstra~\cite{benaloh1994} formalised \emph{receipt-freeness}: the
requirement that a voter cannot construct a convincing proof of their
choice even if they want to. Juels, Catalano and
Jakobsson~\cite{juels2005} strengthened this to
\emph{coercion-resistance}, covering an adversary who can interact with
the voter during voting, demand credentials, or require abstention;
their scheme lets a coerced voter surrender a fake credential
indistinguishable from a real one.

This distinction sets the boundary of the present work. Bribery, as
modelled here, presupposes that the briber can verify compliance well
enough for the transaction to be worth making. A system that is
genuinely receipt-free undermines that presupposition at the root,
which is a different and stronger defence than the detection-based one
modelled in this dissertation. The \emph{Limitations} chapter of the full dissertation returns to this.

\subsection{Blockchain-based voting proposals}
\label{sec:lit-bcvoting}

Nakamoto's chained proof-of-work construction~\cite{nakamoto2008}
supplied a mechanism with obvious surface appeal for elections: an
append-only, hash-linked, replicated log in which retrospective edits
invalidate every subsequent hash and are therefore detectable by anyone
holding a copy.

The most rigorous academic instantiation is McCorry, Shahandashti and
Hao's Open Vote Network~\cite{mccorry2017}, a boardroom-voting protocol
implemented as an Ethereum smart contract, which achieves maximum voter
privacy in the sense that a vote can only be revealed by full collusion
of all remaining participants, and which is self-tallying: no trusted
authority computes the result. It is a genuine cryptographic
contribution rather than a ledger with ballots written into it. Its
scope, however, is explicitly the boardroom --- tens of voters, all
known to one another --- and its cost and collusion properties do not
transfer to a public election.

The wider applied literature is large and, by the assessment of its own
systematic reviewers, uneven. Ohize et al.~\cite{ohize2025} survey
architectures and trends across the field and identify scalability,
post-quantum security, voter authentication, regulatory frameworks and
cross-chain interoperability as the persistent open problems. Rahul,
Gulia and Gill~\cite{rahul2025} review 55 journal articles and reach a
similar conclusion, noting that security, transparency and trust remain
unresolved worldwide rather than settled by any existing proposal. Two
observations from these reviews shape the present work. First,
\emph{voter authentication} is repeatedly named as an open problem,
which is precisely the Sybil vector modelled here. Second, the surveyed
proposals are overwhelmingly evaluated on functionality, throughput and
cryptographic properties rather than on quantified resistance to a
budgeted attacker --- the gap this dissertation addresses.

\subsection{The critical response: what happens on deployment}
\label{sec:lit-critique}

Against the proposal literature stands a smaller and considerably more
sobering body of empirical analysis.

Springall et al.~\cite{springall2014} examined Estonia's I-voting
system, the most established national internet-voting deployment, by
observing operations directly and reconstructing the system in the
laboratory. They reported serious procedural lapses and architectural
weaknesses sufficient, in their assessment, for a state-level attacker
to alter results undetectably, and recommended discontinuing the
system. The finding relevant here is not the specific vulnerabilities
but the location of them: the client platform and the operational
procedures surrounding the count, not the cryptography.

Specter, Koppel and Weitzner~\cite{specter2020} reverse-engineered
Voatz, the first internet-voting application used in U.S. federal
elections and marketed on its use of a blockchain. They found flaws
allowing an adversary to alter, stop or expose a voter's ballot, a
side-channel through which a passive network observer could recover the
secret ballot, and a dependence on third-party services with attendant
privacy exposure; West Virginia and a Washington county withdrew the
application from 2020 primaries following the analysis. The
architectural point is the important one: the blockchain sat behind an
application whose compromise made the ledger irrelevant, because votes
were corrupted before they reached it.

Park, Specter, Narula and Rivest~\cite{park2021} generalise these
observations into a direct argument that blockchain voting is worse
than internet voting rather than better. Their case is that blockchains
address a problem elections do not have --- disagreement about which
transactions occurred among mutually distrusting parties --- while
introducing new failure modes, and that the genuinely hard problems of
elections, namely software independence, secret ballots, voter
authentication and client-side integrity, are untouched. They further
note that a blockchain vote is a \emph{transparent} record, which is in
tension with ballot secrecy rather than supportive of it.

This dissertation adopts that critique as a premise to be tested
quantitatively rather than as a conclusion to be restated. The
experimental design deliberately grants the ledger everything its
proponents claim --- integrity is verified continuously and, as the \emph{Findings} chapter reports, was never violated in any trial ---
and then measures how many elections are stolen anyway.

\subsection{Bribery and vote buying as a formal problem}
\label{sec:lit-bribery}

Three distinct traditions model vote buying, and this dissertation
borrows from all three.

\textbf{Computational.} Faliszewski, Hemaspaandra and
Hemaspaandra~\cite{faliszewski2009} formalise bribery as a decision
problem: given per-voter prices and a budget, can the briber make a
preferred candidate win? They classify the complexity across voting
rules, establishing that some systems make optimal bribery
NP-hard --- complexity as a defensive barrier. The framing is
worst-case, combinatorial and deterministic: a bribed voter votes as
paid.

\textbf{Economic.} Becker's analysis of crime as a decision under
risk~\cite{becker1968} models an offender weighing gain against the
probability of detection multiplied by the penalty, and derives the
policy conclusion that deterrence can be achieved either by raising
detection or by raising penalties. Equation~(\ref{eq:m1}) in this
dissertation is exactly that calculation applied to a bribed voter, and
the critical-detection threshold developed in this work is the
policy question in Becker's sense: how much detection is enough?

\textbf{Blockchain-specific.} Daian, Kell, Bentov and
Juels~\cite{daian2018} describe on-chain vote buying and the
\emph{Dark DAO}: a smart contract that pays for votes and enforces the
payment trustlessly, using trusted execution to verify compliance
without revealing the buyer. Their argument deserves emphasis in the
present context because it inverts the usual claim. A blockchain does
not merely fail to prevent vote buying; by making votes cryptographically
attributable and payments programmable and escrowed, it \emph{improves}
the market for votes, removing the counterparty risk that limits bribery
in traditional settings. This is a technical-report and research-blog
contribution rather than a peer-reviewed paper, and is cited as such,
but it is widely referenced and directly motivates treating bribery,
not tampering, as the primary threat to a ledger-based election.

The gap across all three: the computational tradition treats bribery
deterministically, the economic tradition treats the individual decision
without an electorate-level outcome model, and the blockchain-specific
work is a mechanism proposal rather than a quantified risk model. None
produces the quantity a defender needs --- the probability that a given
budget overturns a given margin against a given detection capability.

\subsection{Sybil attacks and identity}
\label{sec:lit-sybil}

Douceur~\cite{douceur2002} established that, absent a trusted
certification authority, a single entity in a distributed system can
present arbitrarily many distinct identities, and that no purely local
resource test reliably distinguishes them. For an election the
consequence is immediate and fatal: one-person-one-vote is the entire
premise, and a Sybil attack attacks the premise rather than the
mechanism.

Responses divide into those that reintroduce a trusted authority and
those that attempt to establish uniqueness without one. Borge et
al.~\cite{borge2017} propose proof-of-personhood, binding digital
identities to physical presence at pseudonym parties, explicitly to
recover one-per-person semantics in permissionless systems. The
approach illustrates both the necessity and the difficulty: physical
attendance is a strong Sybil barrier and an equally strong barrier to
accessibility, which is what electronic voting exists to improve.

This dissertation does not model any specific identity mechanism. It
parameterises the outcome instead --- a forged registration passes with
probability $q_s$ and is caught with probability $p_d$ --- which is
weaker than assuming a trusted authority and is deliberately
conservative. It also produces the finding, reported in the \emph{Findings} chapter, that identity verification rather than
ledger integrity is where the residual risk concentrates, consistent
with voter authentication being repeatedly flagged as an open problem in
the systematic reviews~\cite{ohize2025,rahul2025}.

\subsection{Ballot stuffing and statistical election forensics}
\label{sec:lit-forensics}

Ballot stuffing --- injecting fraudulent ballots not corresponding to
any voter --- is the least cryptographically interesting of the three
vectors and, historically, among the most common. The relevant
literature is statistical rather than cryptographic.

Klimek, Yegorov, Hanel and Thurner~\cite{klimek2012} show that
systematic irregularities leave detectable signatures in the joint
distribution of turnout and vote share across polling units: incremental
stuffing and extreme fraud produce characteristic distortions absent
from clean elections, and the authors identify these fingerprints in
several national elections. The methodological lesson carried into this
dissertation is that fraud is bounded by \emph{plausibility} as much as
by opportunity. An attacker who injects ballots without limit produces a
turnout figure that requires no probabilistic detection to notice. This
is the justification for the capacity ceilings imposed in the coalition
model (described in the methodology below), which are otherwise a
modelling choice without grounding: they stand in for the anomaly
detection that election forensics performs.

\subsection{Probabilistic and game-theoretic security analysis}
\label{sec:lit-probgame}

The final strand supplies the analytical machinery.

\textbf{Concentration inequalities.} Chernoff's
bound~\cite{chernoff1952}, in the multiplicative form given by
Mitzenmacher and Upfal~\cite{mitzenmacher2017}, bounds the probability
that a sum of independent Bernoulli variables deviates from its mean,
and does so exponentially in the deviation. For this dissertation its
value is specific: it converts an empirical statement into a provable
one. A margin certified safe by the bound is safe irrespective of how
many simulation runs were performed, which is what makes the certified
margin defined in this work a security claim rather than an
observation.

\textbf{Rational attackers.} Eyal and Sirer~\cite{eyal2014}
demonstrated with selfish mining that a blockchain protocol assumed
secure under an honest-majority model can be broken by a rational
minority that deviates profitably --- security depending on incentives,
not only on cryptography or on hash-power share. Liu et
al.~\cite{liu2019} survey the resulting body of game-theoretic
blockchain analysis, covering mining strategy, security games and
incentive design.

Two points transfer directly. First, the relevant question about an
attacker is what they would \emph{choose} to do, not what they are
capable of doing: the \emph{Findings} chapter shows that a
budget-exhausting attacker performs strictly worse than a rational one,
so a defence tuned against maximum expenditure is tuned against the
wrong opponent. Second, deterrence is a legitimate security goal
distinct from prevention --- making an attack unprofitable is cheaper
than making it impossible, and the gap between those two thresholds is
quantified in this dissertation.

\subsection{Synthesis and research gap}
\label{sec:lit-synthesis}

\tabref{tab:gap} summarises what each strand supplies and what it
leaves open.

\begin{table}[htbp]
\centering
\caption{What the literature supplies, and what it leaves open.}
\label{tab:gap}
\small
\begin{tabular}{@{}p{3.5cm}p{5.1cm}p{5.1cm}@{}}
\toprule
\textbf{Strand} & \textbf{Supplies} & \textbf{Leaves open} \\
\midrule
Blockchain voting proposals
  \cite{mccorry2017,ohize2025,rahul2025}
& Architectures, cryptographic protocols, throughput evaluation
& No quantified success probability against a budgeted attacker \\
\addlinespace
Empirical critiques
  \cite{springall2014,specter2020,park2021}
& Evidence that deployments fail upstream of the ledger
& Qualitative argument; no model of how much manipulation gets through \\
\addlinespace
Bribery theory
  \cite{faliszewski2009,becker1968,daian2018}
& Complexity classification; the rational-offender calculus; on-chain vote-buying mechanisms
& Deterministic or individual-level; no electorate-scale outcome distribution \\
\addlinespace
Sybil and identity
  \cite{douceur2002,borge2017}
& Impossibility without trusted certification; countermeasure designs
& No probabilistic model of partial identity verification \\
\addlinespace
Election forensics
  \cite{klimek2012}
& Statistical signatures of stuffing; plausibility as a real constraint
& Detection after the fact, not a design threshold beforehand \\
\addlinespace
Probabilistic and game-theoretic tools
  \cite{chernoff1952,mitzenmacher2017,eyal2014,liu2019}
& Concentration bounds; rational-deviation analysis
& Not yet applied to e-voting manipulation as an integrated model \\
\bottomrule
\end{tabular}
\end{table}

Three gaps follow, and they define the contribution of this
dissertation.

\textbf{Gap 1: no quantified attacker model for e-voting
countermeasures.} The proposal literature evaluates systems on
cryptographic properties and performance; the critical literature
establishes qualitatively that deployments fail. Neither produces the
probability that a specific budget overturns a specific margin against
a specific detection capability. The \emph{Threat Model} and \emph{Findings} chapters supply exactly that, over
\TotalTrials{} matched randomised trials per vector.

\textbf{Gap 2: probabilities are reported where thresholds are
needed.} A defender cannot act on a success-probability curve. The
inverse question --- what detection probability, what budget ceiling,
what margin makes this attack non-viable --- is not asked in the
e-voting literature, although Becker's economics of
deterrence~\cite{becker1968} and the game-theoretic blockchain
work~\cite{eyal2014,liu2019} both imply it should be.
The \emph{Critical Thresholds} chapter formulates seven such thresholds and solves them.

\textbf{Gap 3: attack vectors are evaluated in isolation.} Every source
reviewed treats bribery, Sybil attacks and ballot stuffing separately.
No reviewed work models an attacker allocating one budget across
several vectors simultaneously. The \emph{Coalition Attacks} chapter shows this omission is not conservative: two of the three vectors are individually
infeasible in every benchmark scenario, yet their combination succeeds
decisively in four of five.
